\chapter{Safety first}
\label{ch:u-safety}

The X30 walks among energized high-voltage equipment. Read this chapter before you use the robot
for the first time, and follow it on every shift.

\section{The substation's rules come first}

\begin{itemize}
  \item The robot works under the substation's safety rules and permits. Only people trained on
    this manual and authorized by the substation operate it; one of them is responsible for the
    robot during every patrol, scheduled ones included.
  \item \textbf{Safe distances.} The robot and everything on it stay outside the minimum approach
    distance of every energized part. Teach routes only along walkways and roads people may use
    while the equipment is energized: the robot does not know where live parts are, the route is
    what keeps it away from them.
  \item \textbf{Switching and maintenance.} No patrol through a bay under switching or
    maintenance. Stop the schedule in NETRA and keep the robot on its dock while work goes on.
  \item \textbf{Weather.} No patrols in rain, lightning or standing water, nor on ground the
    robot was not taught on.
\end{itemize}

\section{Four ways to stop the robot}

\begin{center}
\begin{tabularx}{\linewidth}{@{}>{\bfseries}p{0.3\linewidth}X@{}}
  \toprule
  Stop & What happens \\
  \midrule
  Hard emergency stop (tail) & cuts the joints at once: the robot drops. Only when it must stop
    whatever it costs \\
  H12: key \key{B} or STOP & software e-stop: the robot stands, the patrol pauses; release by
    holding STOP for a second \\
  NETRA: \ui{E-stop} & the same e-stop from the dashboard \\
  Deep Robotics' handheld & takes over the robot (manual mode), past the whole stack \\
  \bottomrule
\end{tabularx}
\end{center}

\begin{dangerbox}[The software stops are not certified]
The e-stop of the H12 and of NETRA, and the robot's protective field, are software. They are a
second line, never the only one: keep the hard emergency stop and Deep Robotics' handheld within
reach of someone responsible whenever people are near the robot.
\end{dangerbox}

\section{What the robot watches, and what it does not}

\begin{itemize}
  \item \textbf{It watches} the space around its body with four LiDARs: it slows down and stops
    before obstacles in its way (the protective field), and stops when it cannot see.
  \item \textbf{It does not see} holes, steps down or the edge of a platform; objects lower than
    about 10~cm close to its legs; thin wires. Keep routes on even, solid ground.
  \item \textbf{It does not know} where the live parts are. The route keeps it away from them.
  \item \textbf{It stops itself} when it loses its position, when its sensors fail, or when its
    battery runs low, and says why on the H12 and in NETRA.
\end{itemize}
